\section{计算机如何"看"图像}

\subsection{图像的数字表示}

对于计算机而言，图像仅仅是数字。一幅灰度图像是一个二维数组，其中每个元素（像素）是 $[0, 255]$ 范围内的一个整数，表示该像素的亮度。一幅彩色图像则是一个三维数组，尺寸为 $H \times W \times 3$，其中第三个维度对应 RGB 三个颜色通道。

\begin{figure}[H]
\centering
\includegraphics[width=0.9\textwidth]{slides-images/slide-12.jpg}
\caption{图像就是数字：灰度图像是一个二维矩阵，彩色图像（如 1080$\times$1080）则是三个这样的矩阵堆叠（RGB三通道）。\protect\footnotemark}
\end{figure}
\footnotetext{Slides 第12页。}

\begin{importantbox}{图像的数学表示}
\begin{itemize}
    \item \textbf{灰度图像}：$\mathbf{I} \in \mathbb{R}^{H \times W}$，每个像素值 $I_{i,j} \in [0, 255]$
    \item \textbf{彩色图像}：$\mathbf{I} \in \mathbb{R}^{H \times W \times 3}$，三个通道分别对应红（R）、绿（G）、蓝（B）
    \item 例如，一张 $1080 \times 1080$ 的 RGB 图像包含 $1080 \times 1080 \times 3 = 3,499,200$ 个数字
\end{itemize}
\end{importantbox}

Amini 指出，从表示的角度看，图像实际上比语言更容易处理，因为图像天然就是数字数组，而语言则需要先进行词嵌入（embedding）等转换。

\subsection{计算机视觉的两类基本任务}

计算机视觉中最基本的两类任务是\textbf{回归}（Regression）和\textbf{分类}（Classification）：

\begin{figure}[H]
\centering
\includegraphics[width=0.85\textwidth]{slides-images/slide-13.jpg}
\caption{计算机视觉的基本任务：分类（输出离散类别标签及其概率）和回归（输出连续值）。\protect\footnotemark}
\end{figure}
\footnotetext{Slides 第13页。}

\begin{itemize}
    \item \textbf{回归}：输出变量取连续值，例如预测方向盘的转角。
    \item \textbf{分类}：输出变量取离散类别标签，例如判断图片中的人物是哪位美国总统。分类任务可以输出属于每个类别的概率分布。
\end{itemize}

无论是哪种任务，核心问题都是\textbf{特征检测}（Feature Detection）——理解是什么特征使得不同类别之间存在区分。

\subsection{本章小结}

计算机通过像素值矩阵来表示图像。彩色图像是三维张量（高度 $\times$ 宽度 $\times$ 通道数）。计算机视觉的两类基本任务——分类和回归——都依赖于从图像中有效提取特征的能力。

\newpage
